% https://codeforces.com/contest/2227

\documentclass[oneside,a4paper,14pt]{extarticle} %размер шрифта 14
\usepackage[T1,T2A]{fontenc}
\usepackage[
	a4paper,
	letterpaper,
	top=2cm,
	bottom=2cm,
	left=2.5cm,
	right=1.5cm
]{geometry}
\usepackage[utf8]{inputenc} %кодировка текста
\usepackage[russian]{babel} %поддержка русского языка
\usepackage{textcomp} %текстовые символы
\usepackage{indentfirst} %корректировка отступов
\usepackage{listings} %для кода !!
\usepackage{float} % для "впихивания" изображений в странницу
\usepackage{xcolor}
\usepackage{enumitem} %для нумерации буквами
\usepackage{graphicx} %работа с изображениям
\usepackage{mwe} % for blindtext and example-image-a in example
\usepackage{wrapfig}
\usepackage{caption}
\usepackage{amsmath}  % для формул и символов
\usepackage{amsfonts}
\usepackage{amsthm}
\usepackage{graphicx}
\usepackage[all]{xy}
\usepackage[breaklinks]{hyperref}
\usepackage{xurl}

%%размеркегляузаголоковразделов
\usepackage{titlesec}
\titleformat{\section}
{\normalsize\bfseries}
{\thesection} {1em}{}
\titleformat{\subsection}
{\normalsize\bfseries}
{\thesubsection} {1em}{}
\titleformat{\subsubsection}
{\normalsize\bfseries}
{\thesubsection} {1em}{}

\renewcommand\baselinestretch{1.33}\normalsize % межстрочный интервал
\setlength{\parindent}{1.25cm}
\usepackage{indentfirst}

\lstset{
  language=C,
  basicstyle=\ttfamily\small,
  keywordstyle=\color{blue},
  commentstyle=\color{green!60!black},
  stringstyle=\color{red},
  numbers=left,
  numberstyle=\tiny\color{gray},
  frame=single,
  breaklines=true,
  showstringspaces=false,
  tabsize=2,
  inputencoding=utf8,
  extendedchars=true
}

\begin{document}
	\newpage\thispagestyle{empty}
	\begin{center}
		МИНИСТЕРСТВО НАУКИ И ВЫСШЕГО ОБРАЗОВАНИЯ\\
			РОССИЙСКОЙ ФЕДЕРАЦИИ
			ФЕДЕРАЛЬНОЕ ГОСУДАРСТВЕННОЕ БЮДЖЕТНОЕ\\
			ОБРАЗОВАТЕЛЬНОЕ
			УЧРЕЖДЕНИЕ ВЫСШЕГО ОБРАЗОВАНИЯ\\
			«ВЯТСКИЙ ГОСУДАРСТВЕННЫЙ УНИВЕРСИТЕТ»\\
			Институт математики и информационных систем\\
			Факультет автоматики и вычислительной техники\\
			Кафедра электронных вычислительных машин
	\end{center}
	\vspace{20mm}

	\begin{center}
		Отчёт по лабораторной работе №4\\
		по дисциплине\\
		<<Алгоритмы и структуры данных>>\\
	\end{center}
	\vspace{48mm}

	\noindent
	Выполнил студент гр. ИВТб-1303-06-00 \hfill \rule[-0,5mm]{30mm}{0.15mm}\,/Гортоломей И.К./
	\newline
	Проверил заведующий кафедрой ЭВМ \hfill \rule[-0,5mm]{30mm}{0.15mm}\,/Долженкова М.Л./

	\vfill
	\begin{center}
		Киров\\
		2026
	\end{center}

	\newpage
	% \thispagestyle{empty}

	\section*{Цель работы:}
	Отточить навыки решения задач в условиях ограниченного времени.

  \section*{Вводные}
  Для выполнения данной лабораторной работы был выбран Codeforces Round 1096 (Div. 3). 
  \url{https://codeforces.com/contest/2227}

  \section*{Описание процесса решения}
  Было выбрано и запущено виртуальное соревнование.
  Первые 40 минут соревнования я пропустил, так как пришлось отвлечься и написать контрольную по истории.
  Далее я открыл задачу A - Кошари \url{https://codeforces.com/contest/2227/problem/A},
  в данной задаче требуется дать ответ yes/no, можно ли добраться из точки (0, 0) в (x, y)
  имея возможность прибавить 2 к одной из координат за шаг или прибавить 1 к одной из координат один раз.
  Если мы можем прибавить по 2 сколько угодно раз, то значит мы передвигаемся по чётным,
  а также имея возможность один раз прибавить 1 мы можем сдвинуться на нечётное по одной из координат,
  следовательно, сумма остатков от деления x и y на 2 должна быть меньше двух. $x\%2 + y\%2 < 2$.

  \begin{lstlisting}[caption={Решение задачи A - Кошари}]
#include <stdio.h>

int main() {
	int t; scanf("%d", &t);
	for (; t > 0; t--) {
		int x, y; scanf("%d %d", &x, &y);
		printf(x%2 + y%2 > 1 ? "no\n" : "yes\n");
	}; return 0;
}
  \end{lstlisting}
  \url{https://codeforces.com/contest/2227/submission/374474673}
  \newpage

  После быстрого решения первой задачи, я приступл к решению следующей задачи.
  В задаче B. Party Monster (\url{https://codeforces.com/contest/2227/problem/B}) требуется вывести yes/no
  в зависимости от того, можно ли составить из входящих байт, скобок (открывающая <<(>> и закрывающая <<)>>)
  правильную структуру (например <<()()>> или <<(())>>). Первым делом я решил проверить, что пройдёт ли код
  который просто проверяет кол-во открывающих и закрывающих скобок, ведь если их разное кол-во, то не получится
  составить из них правильную структуру. Быстро написал код который считает скобки $a == b$ и это решение подошло.

  \begin{lstlisting}[caption={Решение задачи B. Party Monster}]
#include <stdio.h>

int main() {
	int t; scanf("%d", &t);
	for (; t > 0; t--) {
		int n; scanf("%d\n", &n);
		int a = 0;
		int b = 0;
		for (; n > 0; n--) {
			char x = getchar();
			if (x == '(') a++;
			if (x == ')') b++;
		}
		printf(a != b ? "no\n" : "yes\n");
	}
	return 0;
}
  \end{lstlisting}
  \url{https://codeforces.com/contest/2227/submission/374475638} \\

  После решения первых двух заданий осталось примерно 1ч 30м до конца, так что я приступил к следующей задаче.
  \newpage

  Следующая задача C - Снегопад (\url{https://codeforces.com/contest/2227/problem/C}) не показалась простой сразу.
  Минут 10 не мог понять, что требует задача и почему получаю неверный ответ. Оказалось, что я просто забыл развести числа,
  которые делятся на 2 и/или 3 и могут образовывать последовательности делящиеся на 6. Сделал вывод по массиву, чтобы сначала 
  сначала выводить числа, которые делятся на 6, потом числа которые делятся на 2 (но не делятся на 6), за ними <<нейтральные>> числа,
  которые не делятся ни на 2, ни на 3, а в конце числа которые делятся на 3 (но не делятся на 6). В итоге я получаю массив, в котором
  кол-во подмасивов, делящихся на 6 минимальное. Обработка занимает O(4n) -> O(n).

  \begin{lstlisting}[caption={Решение задачи C - Снегопад}]
#include <stdio.h>

int main() {
	int t;
	scanf("%d", &t);
	
	while (t--) {
		int n;
		scanf("%d", &n);
		
		int a[n];
		for (int i = 0; i < n; i++) {
			scanf("%d", &a[i]);
		}

		for (int i = 0; i < n; i++) {
			if (a[i] % 6 == 0) {
				printf("%d ", a[i]);
			}
		}
		
		for (int i = 0; i < n; i++) {
			if (a[i] % 6 != 0 && a[i] % 2 == 0) {
				printf("%d ", a[i]);
			}
		}

		for (int i = 0; i < n; i++) {
			if (a[i] % 2 != 0 && a[i] % 3 != 0) {
				printf("%d ", a[i]);
			}
		}

		for (int i = 0; i < n; i++) {
			if (a[i] % 6 != 0 && a[i] % 3 == 0) {
				printf("%d ", a[i]);
			}
		}
		
		printf("\n");
	}
	
	return 0;
}
  \end{lstlisting}
  \url{https://codeforces.com/contest/2227/submission/374479599}
  В итоге, после решения этой задачи оставалось около 45 минут до конца соревнования,
  потому что я забыл про делители 2 и 3, и по этому потратил много времени.

  Я посмотрел задачу D. ПалиндроМЕХ (\url{https://codeforces.com/contest/2227/problem/D}), понял что нужно
  как-то перебрать все палиндромы в массиве, выполнить функцию mex для каждого палиндрома (половины палиндрома)
  и найти такой палиндром, у которого функция mex возвращает максимальное число (mex возвращает минимальное отстуствующее число).
  На обдумывание этой задачи ушло 5-10 минут, так что я роешил не пытаться писать алгоритм и перешёл к следующей задаче.

  \newpage

  В задаче E. Всё съехало (\url{https://codeforces.com/contest/2227/problem/E}) нам дают одномерный массив (который как бы 
  представляет из себя матрицу) блоков, и нам нужно сказать сколько блоков сдвинется, когда гравитация поменяется с направления вниз
  на направление вправа. Данная задача решалась бы просто, если бы не одно НО: нам предоставляется возможность убрать один из кубов,
  чтобы другие кубы могли сдвинуться и нам нужно выбрать такой куб, после удаления которого сдвинется максимальное число кубов.
  Я начал писать алгоритм, проверил что подсчёт сдвинувшихся правильный, но не правильно выбирал куб который нужно удалять. Я предполагаю,
  что где-то в массиве могут быть столбцы в которых много кубов, а между ними столбец с 1-2 кубами и при удалении куба в этом столбце
  сдвинется много других кубов. Когда я это понял - время соревнования закончилось и я не успел завершить задачу.

  \section*{Выводы}
  Проанализировав своё участие в соревновании я пришёл к выводу, что мне нужно больше времени для обдумывания поставленных условий,
  подбора подходящего алгоритма и проверки кода на ошибки для получения качественного результата. Много времени уходит на промежуточные
  проверки (компиляция тестов, дебагинг и тд.) потому что мне необходимо <<получать обратную связь от кода>>. \\
  \newline
  \indent
  Для ускорения решения задач можно отказаться от промежуточных проверок и заучить синтаксис одного из удобных языков,
  но <<ускорить мышление>> в моей ситуации наверное не получится, или это придёт с опытом.

\end{document}
